\section{Threats to Validity}
Internal validity: (1) The Codec ground truth failed on the first run (100\% NO\_COVERAGE, caused by the absence of JUnit 4 - JUnit 3 tests could not be executed by PIT): had this gone unnoticed, the entire Codec dataset would have been noise, dragging all comparisons toward zero. Mitigation: an automatic per-project guard (which fails when a project has 0 killed mutants or $>50\%$ NO\_COVERAGE) caught the error; PIT was rerun with the corrected classpath and verified on one version before running all 10; CSV and JacksonCore were confirmed to have no byte-for-byte changes; the corrupted version was saved as proof. (2) The raw GPT scores (1-5) create ties in the selection: the order in which tied mutants are selected can affect the results. Mitigation: ties are broken by mutant\_id (a fixed rule, independent of the outcome); selection is done by group so that ties do not accumulate within a single project. (3) The random baseline uses a single seed: one seed could be lucky or unlucky. Mitigation: seed 42 was pre-registered in the proposal; a Monte Carlo analysis of 5,000 draws placed seed 42 at the 52nd percentile of the random distribution - a typical, unbiased representative. (4) Non-determinism of the LLM: temperature=0, a pinned model version, and per-batch checkpointing were used; however, the vendor does not guarantee absolute determinism across calls - this residual risk is noted but not fully eliminated.

External validity: The scope covers 3 Java/Defects4J projects, 20 classes, and 37 groups - the conclusion is limited to this scope and is not generalized to all Java projects, other languages, or other LLMs. Mitigation: the three projects belong to three different domains (encoding/CSV/JSON) rather than the same utility family; the report clearly states that the number of classes (not the number of projects) drives the sample size n - a lesson learned from the fact that our 3 projects have noticeably fewer groups than designs with 2 projects but many more classes.

Construct validity: (1) MS = Killed/Selected is computed while still including NO\_COVERAGE mutants in the selected sample - it would be reasonable to exclude them from the denominator. Mitigation: measured - the NO\_COVERAGE ratios in the two branches are close (GPT 20.7\% vs. random 21.8-23.5\%), and on the covered-only subset, GPT still performs better (0.864 vs. 0.810); the conclusion remains unchanged under this alternative definition. (2) The "usefulness" that the prompt asks GPT to score is not entirely consistent with the "killability" that MS measures - a mutant that is hard to kill (survives) may still be useful in that it points to a test suite weakness. We note this as a construct limitation; future studies may instead measure fault-revelation, as in \cite{jimenez2018naturalness}.

Conclusion validity: (1) N=37 groups gives low statistical power for group-weighted testing; 11/37 pairs are tied, which degenerates the median criterion. Mitigation: an amendment removed the median criterion for structural reasons (approved by the instructor and documented before finalizing the conclusion); mutant-weighted and Monte Carlo results are reported in parallel; the conclusion is framed as "partial support" rather than forced into a binary outcome. (2) The CSV subgroup has n=10 - right at the pre-registered threshold (n<10 → descriptive only); subgroup results should be read only as a reference. (3) Multiple analyses were run on the same data (risk of multiple testing). Mitigation: the pre-registered analysis is reported as-is and first; robustness analyses are clearly marked as post-hoc and have methodological justifications independent of the results.